\documentclass[tikz,border=6pt]{standalone}
\usepackage{peabra}
\begin{document}
\begin{tikzpicture}[x=1cm,y=1cm]
  \tikzset{f/.style={doboz=33mm, font=\scriptsize, minimum height=12mm, align=center},
           g/.style={zolddoboz=33mm, font=\scriptsize, minimum height=12mm, align=center}}
  % sávok
  \fill[peHalvany!55] (-2.3,0.95) rectangle (17.1,-0.95);
  \fill[peZoldH!55] (-2.3,-2.65) rectangle (17.1,-4.55);
  \node[alcim, rotate=90, align=center] at (-1.95,0) {táblázat};
  \node[alcim, rotate=90, align=center] at (-1.95,-3.6) {notebook};
  % táblázatkezelő
  \node[f] (A) at (0,0) {\textbf{1. Eszközlap}\\lépések, kidolgozott példa};
  \node[f] (B) at (3.8,0) {\textbf{2. Excel-sablon}\\a kitöltött példa áttekintése};
  \node[f] (C) at (7.6,0) {\textbf{3. Saját adatok}\\a sárga cellákba};
  \node[f] (D) at (11.4,0) {\textbf{4. Ellenőrzés}\\képletek, figyelmeztetések};
  % notebook
  \node[g] (N1) at (3.8,-3.6) {\textbf{5. Notebook Colabban}\\a beépített példa futtatása};
  \node[g] (N2) at (7.6,-3.6) {\textbf{6. Saját sablon betöltése}\\\texttt{SAJAT\_SABLON = True}};
  \node[g] (N3) at (11.4,-3.6) {\textbf{7. További elemzés}\\szimuláció, gráf, statisztika};
  % eredmény
  \node[sotetdoboz=30mm, font=\scriptsize, minimum height=12mm, align=center] (R) at (15.4,-1.8) {\textbf{8. Eredmények}\\a projektjelentésbe};
  \foreach \a/\b in {A/B,B/C,C/D,N1/N2,N2/N3} \draw[nyil] (\a) -- (\b);
  \draw[nyil] (D.east) -| node[felirat, above, pos=0.25] {elég a táblázat} (R.north);
  \draw[nyil] (N3.east) -| (R.south);
  \draw[nyil] (D.south) -- node[felirat, right, align=left] {további\\elemzés kell} ++(0,-1.15) -| (N2.north);
  \draw[nyil, dash pattern=on 3pt off 2pt] (B.south) -- node[felirat, left, align=right] {a módszer\\kipróbálása} (N1.north);
\end{tikzpicture}
\end{document}
